% Part:sets-functions-relations
% Chapter: sets
% Section: reduction

\documentclass[../../../include/open-logic-section]{subfiles}

\begin{document}

\olfileid[fa-IR]{sfr}{siz}{red}

\olsection{کاهش}

\begin{editorial}
  این بخش ناشمارابودن را از راه کاهش ثابت می‌کند و با نتایجِ
  \olref[nen]{sec} مطابقت دارد. روایتی جایگزین و اندکی موجزتر که با
  نتایجِ \olref[nen-alt]{sec} مطابقت دارد، در \olref[red-alt]{sec}
  آمده است.
\end{editorial}

با استدلالی قطری نشان دادیم که $\Pow{\PosInt}$ !!{nonenumerable} است.
از پیش نیز برهانی داشتیم که $\Bin^\omega$، یعنی مجموعهٔ همهٔ
دنباله‌های نامتناهیِ $0$ها و $1$ها، !!{nonenumerable} است. راه دیگری
برای اثبات اینکه $\Pow{\PosInt}$ !!{nonenumerable} است چنین است: نشان دهید
که \emph{اگر $\Pow{\PosInt}$ !!{enumerable} باشد، آنگاه $\Bin^\omega$
  نیز !!{enumerable} است}. چون می‌دانیم $\Bin^\omega$ !!{enumerable}
نیست، $\Pow{\PosInt}$ نیز نمی‌تواند چنین باشد. این کار را
\emph{کاهش‌دادن} یک مسئله به مسئله‌ای دیگر می‌نامند---در اینجا، مسئلهٔ
برشمردن $\Bin^\omega$ را به مسئلهٔ برشمردن $\Pow{\PosInt}$ کاهش
می‌دهیم. حل مسئلهٔ دوم---یعنی برشماریِ $\Pow{\PosInt}$---راه‌حلی برای
مسئلهٔ نخست---یعنی برشماریِ $\Bin^\omega$---به دست می‌دهد.

چگونه مسئلهٔ برشمردن مجموعهٔ~$B$ را به مسئلهٔ برشمردن مجموعهٔ~$A$
کاهش می‌دهیم؟ روشی ارائه می‌کنیم که برشماریِ~$A$ را به برشماریِ~$B$
تبدیل کند. آسان‌ترین راه آن است که تابعی !!a{surjective} به‌صورت
$f\colon A \to B$ تعریف کنیم. اگر $x_1$، $x_2$، \dots{} مجموعهٔ~$A$ را
برشمارد، آنگاه $f(x_1)$، $f(x_2)$، \dots{} مجموعهٔ~$B$ را برمی‌شمارد.
در مورد ما، در پی تابعی پوشا به‌صورت
$f\colon \Pow{\PosInt} \to \Bin^\omega$ هستیم.

\begin{prob}
نشان دهید اگر تابعی !!{injective} به‌صورت $g\colon B \to A$ وجود داشته
باشد و $B$~!!{nonenumerable} باشد، آنگاه $A$ نیز چنین است. برای این
کار نشان دهید چگونه می‌توان با استفاده از~$g$، برشماریِ~$A$ را به
برشماریِ~$B$ تبدیل کرد.
\end{prob}

\begin{proof}[برهانِ {\olref[nen]{thm:nonenum-pownat}} از راه کاهش]
فرض کنید $\Pow{\PosInt}$ !!{enumerable} باشد و بنابراین برشماری‌ای از
آن به‌صورت $Z_{1}$، $Z_{2}$، $Z_{3}$، \dots وجود داشته باشد.

تابع $f \colon \Pow{\PosInt} \to \Bin^\omega$ را با قراردادنِ
$f(Z)$ برابر دنبالهٔ $s_{k}$ تعریف کنید، به‌گونه‌ای که $s_{k}(n) = 1$ اگر و تنها
اگر $n \in Z$، و در غیر این صورت $s_k(n) = 0$. این امر به‌روشنی تابعی
تعریف می‌کند، زیرا هرگاه $Z \subseteq \PosInt$، هر $n \in \PosInt$ یا
!!a{element}ی از $Z$ است یا نیست. برای نمونه، مجموعهٔ $2\PosInt = \{2,
4, 6, \dots\}$ از اعداد صحیح مثبت زوج به دنبالهٔ $010101\dots$ نگاشته
می‌شود، مجموعهٔ تهی به $0000\dots$ و خود مجموعهٔ $\PosInt$ به
$1111\dots$ نگاشته می‌شود.

این تابع همچنین !!{surjective} است: هر دنباله از $0$ها و $1$ها با
مجموعه‌ای از اعداد صحیح مثبت متناظر است؛ همان مجموعه‌ای که اعضایش
اعداد صحیحِ متناظر با جایگاه‌هایی هستند که دنباله در آن‌ها مقدار~$1$
دارد. دقیق‌تر بگوییم، فرض کنید $s \in \Bin^\omega$. مجموعهٔ $Z
\subseteq \PosInt$ را چنین تعریف کنید:
\[
Z = \Setabs{n \in \PosInt}{s(n) = 1}
\]
آنگاه $f(Z) = s$؛ این را می‌توان با مراجعه به تعریفِ~$f$ بررسی کرد.

اکنون فهرست زیر را در نظر بگیرید:
\[
f(Z_1), f(Z_2), f(Z_3), \dots
\]
چون $f$ !!{surjective} است، هر عضوِ $\Bin^\omega$ باید برای ورودی‌ای
به‌عنوان مقدارِ~$f$ ظاهر شود و بنابراین باید در فهرست بیاید. پس این
فهرست باید همهٔ~$\Bin^\omega$ را برشمارد.

بنابراین اگر $\Pow{\PosInt}$ !!{enumerable} بود، $\Bin^\omega$ نیز
!!{enumerable} می‌بود. اما $\Bin^\omega$ !!{nonenumerable} است
(\olref[nen]{thm:nonenum-bin-omega}). پس $\Pow{\PosInt}$
!!{nonenumerable} است.
\end{proof}

\begin{explain}
ممکن است جهت کاهش به‌آسانی موجب سردرگمی شود. برای نمونه، تابعی
!!a{surjective} مانند $g \colon \Bin^\omega \to B$ \emph{ثابت نمی‌کند}
که $B$ !!{nonenumerable} است. (تابع $g
\colon \Bin^\omega \to \Bin$ را در نظر بگیرید که با $g(s) = s(1)$
تعریف می‌شود و هر دنباله از $0$ها و $1$ها را به نخستین !!{element} آن
می‌نگارد. این تابع !!{surjective} است، زیرا بعضی دنباله‌ها با $0$ و
بعضی با $1$ آغاز می‌شوند؛ اما $\Bin$ متناهی است.) همچنین توجه کنید که
تابع~$f$ باید !!{surjective} باشد، وگرنه استدلال پیش نمی‌رود:
فهرستِ $f(x_1)$، $f(x_2)$، \dots{} در آن صورت تضمین نمی‌کند که همهٔ
!!{element}s را در~$B$ شامل شود. برای نمونه،
\[
h(n) = \underbrace{000\dots0}_{\text{$n$ تا $0$}}
\]
تابعی به‌صورت $h\colon \PosInt \to
\Bin^\omega$ تعریف می‌کند، اما $\PosInt$ !!{enumerable} است.
\end{explain}

\begin{prob}
با استدلالی کاهشی نشان دهید که مجموعهٔ همهٔ \emph{مجموعه‌های} زوج‌های
اعداد صحیح مثبت !!{nonenumerable} است.
\end{prob}

\begin{prob}\label{sfr:siz:red:prob:nat-nat}
  با استدلالی کاهشی نشان دهید مجموعهٔ~$X$ از همهٔ توابع
  $f\colon \Nat \to \Nat$ !!{nonenumerable} است. (راهنمایی: تابعی پوشا
  از $X$ به~$\Bin^\omega$ بدهید.)
\end{prob}

\begin{prob}
با استدلالی کاهشی نشان دهید $\Nat^\omega$، یعنی مجموعهٔ دنباله‌های
نامتناهیِ اعداد طبیعی، !!{nonenumerable} است.
\end{prob}

\begin{prob}
فرض کنید $P$ مجموعهٔ توابع از مجموعهٔ اعداد صحیح مثبت به مجموعهٔ
$\{0\}$ باشد و $Q$ مجموعهٔ توابع \emph{جزئی} از مجموعهٔ اعداد صحیح
مثبت به مجموعهٔ $\{0\}$ باشد. نشان دهید $P$~!!{enumerable} است و
$Q$~نیست. (راهنمایی: مسئلهٔ برشمردن $\Bin^\omega$ را به مسئلهٔ
برشمردن~$Q$ کاهش دهید.)
\end{prob}

\begin{prob}
فرض کنید $S$ مجموعهٔ همهٔ توابع !!{surjective} از مجموعهٔ اعداد صحیح
مثبت به مجموعهٔ \{0,1\} باشد؛ یعنی $S$ از همهٔ توابع
!!{surjective}~$f\colon \PosInt \to \Bin$ تشکیل می‌شود. نشان دهید $S$
!!{nonenumerable} است.
\end{prob}

\begin{prob}
نشان دهید مجموعهٔ~$\Real$ از همهٔ اعداد حقیقی !!{nonenumerable} است.
\end{prob}

\end{document}
